% kernel/engineering_context.tex

\chapter{Engineering Context}
\label{chap:engineering-context}

Engineering information is interpreted within a context. Context
provides the conditions under which an engineering object, state,
observation, realization, or decision obtains its applicable meaning.

\section{Context}

\begin{definition}[Engineering context]
\label{def:engineering-context}
Engineering context is the set of applicable conditions required to
interpret engineering information without altering the identity of the
engineering object to which that information refers.
\end{definition}

Context may include a project, design standard, operational condition,
metric realization, coordinate reference system, observation setting,
construction phase, or decision authority.

\section{Context and Identity}

Context qualifies engineering information but does not automatically
become part of object identity. The same engineering object may
participate in several projects, be represented in several coordinate
systems, or be evaluated under several operational conditions.

A context-specific description must therefore remain distinguishable
from the durable object that it qualifies.

\section{Context and State}

States are meaningful only when their applicable context is known. A
geometric state, measured state, target state, and operational state may
refer to the same engineering object while answering different
engineering questions.

Explicit context prevents values with different meanings from being
combined merely because they share a numerical form.

\section{Context and Authority}

Context also determines which rules, knowledge, and authorities apply.
A proposal admissible within one design context may be inadmissible in
another. An observation valid for one realization context may not be
transferable without qualification.

\section{Canonical Interpretation}

\begin{proposition}
Engineering context qualifies interpretation without replacing
engineering identity.
\end{proposition}

\begin{proposition}
Context-dependent information shall not be treated as context-free
engineering truth.
\end{proposition}

\begin{proposition}
The applicability of states, observations, constraints, and decisions
depends on an explicit engineering context.
\end{proposition}

Engineering context provides the interpretive boundary through which
canonical engineering concepts become applicable to concrete
engineering questions.
